\section{Code Snippets for Reverberation}
\label{app:reverb_code}

This appendix presents the full set of implementation for the Signalsmith's Reverb
described in Section \ref{sec:reverb-discussion}.

\subsection{Multi-channel Comb Filter}

The following implements the delay structure described by equation \eqref{eq:delay}.

\lstinputlisting[
    style=cppstyle,
    linerange={process_delay-process_delay},
    caption={Delay line processing},
    label={lst:process-delay}
]{../include/audio_fx/delay/Delay.hpp}

The listing below (\Cref{prompt:householder_code}) \cite{openai_chatgpt} shows the implementation for both 4-order and 8-order Householder 
transform \eqref{eq:householder-general}.

\lstinputlisting[
    style=cppstyle,
    linerange={householder-householder},
    caption={Householder matrix transformation},
    label={lst:householder}
]{../include/audio_fx/reverb/Hadamard.hpp}

These array of delays and the Householder matrix are implemented in the multi-channel IIR Comb Filter \eqref{eq:feedback-loop} shown below.

\lstinputlisting[
    style=cppstyle,
    linerange={process_feedback-process_feedback},
    caption={Algorithm for the Multi-channel Comb Filter},
    label={lst:process_feedback}
]{../include/audio_fx/reverb/allpass_systems/MultiFeedback.hpp}


\subsection{Diffusion System}

After delaying the signals (Listing \ref{lst:process-delay}), the signals are shuffled and some of them are inverted.
The listings below shows how the array of signals are shuffled and inverted.

\lstinputlisting[
    style=cppstyle,
    % linerange={process_feedback-process_feedback},
    caption={Shuffles the indices of the channels},
    label={lst:indices-shuffle}
]{appendices/code/permutation.cpp}

\lstinputlisting[
    style=cppstyle,
    % linerange={delay_shuffle-delay_shuffle},
    caption={Inversion of amplitudes through natural selection},
    label={lst:polarity_change}
]{appendices/code/polarity.cpp}

And finally, the array of signals pass through Hadamard transform, which is implemented in the listing below,
for both 4-order and 8-order.
\lstinputlisting[
    style=cppstyle,
    linerange={hadamard-hadamard},
    caption={Hadamard matrix transformation},
    label={lst:hadamard}
]{../include/audio_fx/reverb/Hadamard.hpp}

After putting everything together, the listing below shows how diffusion is processed.
\lstinputlisting[
    style=cppstyle,
    linerange={diffusion_process-diffusion_process},
    caption={Complete diffusion process},
    label={lst:diffusion_process}
]{../include/audio_fx/reverb/allpass_systems/Diffusion.hpp}


\subsection{Other Helper Functions}

The listing below duplicates the signal to 4, 8, or sometimes 16. The program should implement it
at the very beginning of the reverberation process.
\lstinputlisting[
    style=cppstyle,
    % linerange={split-split},
    caption={Splitting function},
    label={lst:split}
]{appendices/code/split.cpp}

The listing below reunites the signals to one. The program should implement it
at the very end of the reverberation process.
\lstinputlisting[
    style=cppstyle,
    % linerange={split-split},
    caption={Joining function},
    label={lst:join}
]{appendices/code/join.cpp}

Since the amplitudes are added, it could cause clips and distortions during resonance.
So, sometimes, a gain limiter needs to be implemented.

The listing below generates delays by taking the mean of the Delay time (in milliseconds),
converts it to samples, and then takes a random range of numbers within 30-40\% of the mean.
The quantity depends on the number of channels used.

\lstinputlisting[
    style=cppstyle,
    linerange={generate_delays-generate_delays},
    caption={Generates a random set of delays close to the mean},
    label={lst:generate_delays}
]{../include/audio_fx/reverb/ReverbUtils.hpp}

That particular code is only used in the comb filter. However, the diffusion system uses
a slightly different approach. The function takes the given largest delay time it can generate
and then generates a set of delay samples from 0 to largest. The algorithm for generating,
however, remains the same.

% TODO: Complete process from reverb
\subsection{Complete Reverberation Process}

Using these algorithms mentioned above, everything is put together to create a
complete algorithmic reverberator below.

\lstinputlisting[
    style=cppstyle,
    linerange={complete_process-complete_process},
    caption={The complete reverberation process},
    label={lst:complete_reverb_process}
]{../include/audio_fx/reverb/Reverb.cpp}
